\documentclass[10pt,a4paper]{amsart}
\usepackage[margin=19mm]{geometry}
\usepackage{amsmath,amssymb,mathtools}
\usepackage[scr=rsfs,cal=euler]{mathalfa}
\usepackage{xfrac}
\usepackage[unicode,hidelinks,bookmarks=false]{hyperref}
\newcommand{\ie}{i.e.\ }
\newcommand{\cf}{cf.\ }
\newcommand{\resp}{respectively\ }
\newcommand{\Subsection}{\subsection}
\providecommand{\nobreakdash}{\nobreak\hbox{-}}
\setlength{\emergencystretch}{2em}
\multlinegap0pt
\usepackage{amssymb}
\usepackage{bbm}
\usepackage{dsfont}
\usepackage{caption}
\newtheorem{lemma}{Lemma}
\title{Low-degree Lower bounds for clustering in moderate dimension}
\author{Alexandra Carpentier, Nicolas Verzelen}
\date{Source: arXiv:2602.23023v1 (CC BY 4.0)}
\begin{document}
\maketitle

\noindent\emph{Source and licence.} Low-degree Lower bounds for clustering in moderate dimension. Version arXiv:2602.23023v1 (CC BY 4.0), distributed by arXiv under the Creative Commons Attribution 4.0 International licence, http://creativecommons.org/licenses/by/4.0/, as recorded in the arXiv licence metadata for this exact version and shown on the abstract page https://arxiv.org/abs/2602.23023v1. Full licence notice: \texttt{licenses/arxiv-2602.23023-CC-BY-4.0.txt}.\par\smallskip

\noindent\textit{Editorial scope.} Counted unit: the article's lemma lem:terme\_correctif (source0.tex lines 2312--2317). The article's complete original proof of it is delivered with it (source0.tex lines 2323--2352, one \texttt{proof} environment of 30 lines, which the article places away from the statement). No mathematics is written, repaired or re-derived: the statement and the proof are the article's own lines, in the article's own notation, and no retained line is substituted or re-flowed.  No further source-local context is needed: every cross-reference of every retained line resolves inside the two delivered spans.  Nothing was omitted from any retained span, so the package carries no omission record.  No retained line contains a citation, so the wrapper carries no bibliography and no bibliography entry.  No retained line contains a figure environment, an image-inclusion command or drawing code, so no image is delivered and the package declares no asset dependency.  Editorial formatting only, in addition: an AMS \texttt{amsart} preamble that restates the article's own declarations and macros used by the retained lines, namely the environments \texttt{lemma} on separate counters, together with the standard packages those lines need.  The retained environments print in this document as Lemma 1 in order of appearance, whereas the article prints the same items with the numbering of its own full document.  The complete native source of the article as distributed in the arXiv e-print for this exact version is bundled unchanged as \texttt{sources/195366/source0.tex}.
    \begin{lemma}\label{lem:terme_correctif}
    If $V_{2,np}\neq \emptyset$, we claim that 
    \[
    \left|   \mathbb{E}\left[ \underline{\omega}_{\mathbf{P}}\right]\right| \leq 2^{|V_{2,np}|}\mathbb{E}\left[\omega_{ \mathbf{P}}\right]\ . 
    \]
    \end{lemma}

    \begin{proof}[Proof of Lemma~\ref{lem:terme_correctif}]
    For any node $v$ in $V_{2,np}$, there is a correction term $-(\Delta^2/K) \mathbf{1}_{j^{(*)}_{\underline{e}}= j^{(*)}_{\underline{e}'}}$ where $(*)=(1),(2)$ depending on whether $v$ belongs to $V^{(1)}$ or $V^{(2)}$ and where $\underline{e}$ and $\underline{e}'$ are the two-half-edges incident to $v$. 
    Let us expand all these terms in $\underline{\omega}_{\mathbf{P}}$. For any $T\subseteq V_{2,np}$, we consider the term of $\underline{\omega}_{\mathbf{P}}$ where we consider the correction term 
    $-(\Delta^2/K) \mathbf{1}_{j^{(*)}_{\underline{e}}= j^{(*)}_{\underline{e}'}}$ for all $v\in T$ and, for any $v\notin T$ we consider the term 
    $\mu_{k^{*}(\pi^{(*)}(v)),j^{(*)}_{\underline{e}}} \mu_{k^{*}(\pi^{(*)}(v)),j^{(*)}_{\underline{e}'}}$ where $(*)=1,2$ depending whether $v\in V^{(1)}$ or $V^{(2)}$. In order to compute the corresponding quantity, we build a graph $G_{\Delta,T}=(V_{\Delta},E_{\Delta,T})$ from $G_{\Delta}$ as follows: first, for all $v\in T$, we color (with a color that depends on $v$) all half-edges incident to $v$ that were originally in the same graph $G^{(1)}$ or $G^{(2)}$. Then, we build $G_{\Delta,T}$ by removing all possible cycles formed by colored half-edges and removing all half-edges that formed an open path by connecting the two (non-colored) extremities of these open paths. Writing $|\mathrm{Cyc}(T)|$ the number of cycles that have been pruned when going from $G_{\Delta}$ to $G_{\Delta,T}$. Then, we have 
    \begin{eqnarray*}
        \underline{\omega}_{\mathbf{P}} &:= \sum_{T\subseteq V_{2,np}} (-1)^{|T|} d^{|\mathrm{Cyc}|} d^{|\mathrm{Cyc}(T)|} \frac{\Delta^{2|T|}}{K^{T}}\prod_{e\in E_{\Delta,T}} \langle \mu_{k^*(l(e))},  \mu_{k^*(r(e))}\rangle\ .
    \end{eqnarray*}
    Integrating with respect to $k^*(.)$ and using the orthogonality  of the $\mu_k$'s, we arrive at 
    \begin{eqnarray}\nonumber
    \mathbb{E}\left[\underline{\omega}_{\mathbf{P}}\right]&=& \sum_{T\subseteq V_{2,np}} (-1)^{|T|}  d^{|\mathrm{Cyc}(T)|} \frac{\Delta^{2|T|}}{K^{|T|}} \frac{\Delta^{2|E_{\Delta,T}|}}{K^{|V_{\Delta}|-|CC(G_{\Delta,T})|}} \\
    & = &   \frac{d^{|\mathrm{Cyc}(T)|} \Delta^{2|E_{\Delta}|}}{K^{|V_{\Delta}|-|CC(G_{\Delta})|}}\sum_{T\subseteq S} (-1)^{|T|}  \frac{d^{|\mathrm{Cyc}(T)|}}{K^{|T|+|CC(G_{\Delta})| -|CC(G_{\Delta,T})|}} \enspace . \label{eq:upper:corrective}
    \end{eqnarray}
    To conclude, we use that $d\leq K$ and we claim that  
    \begin{equation}\label{eq:upper_T}
    |T|+ |CC(G_{\Delta})| -|CC(G_{\Delta,T})|- \mathrm{Cyc}(T)\geq 0\enspace . 
    \end{equation}
    We first conclude the proof and then we show this claim. We deduce from~\eqref{eq:upper:corrective} and~\eqref{eq:upper_T} that 
    \[
    \left|\mathbb{E}\left[\underline{\omega}_{\mathbf{P}}\right]\right|\leq \frac{d^{|\mathrm{Cyc}(T)|} \Delta^{2|E_{\Delta}|}}{K^{|V_{\Delta}|-|CC(G_{\Delta})|}} 2^{|V_{2,np}|}\ , 
    \]
    which is the desired bound. 
    
    Let us now show~\eqref{eq:upper_T}  by induction on $T$. This is obviously true for $T=\emptyset$. Suppose this is true for $T\setminus\{v\}$. When we add $v$, we prune two more half-edges to go from $G_{\Delta,T\setminus \{v\}}$ to $G_{\Delta,T}$. It these two half-edges corresponded to a self-edge of $v$, then the number of cycles increases by one, we number of connected components is the same:  
    \[
    |T|- |CC(G_{\Delta})| -|CC(G_{\Delta,T})|- |\mathrm{Cyc}(T)| = |T\setminus\{v\}|- |CC(G_{\Delta})| -|CC(G_{\Delta,T\setminus\{v\}})|- |\mathrm{Cyc}(T\setminus\{v\})| \geq 0\ . 
    \]
    It these two half-edges do not correspond to a self-edge, then $\mathrm{Cyc}(T\setminus\{v\})= \mathrm{Cyc}(T)$. Besides, by going from $G_{\Delta,T\setminus\{v\}}$ to $G_{\Delta,T}$, we have replaced two edges incident to $v$ by an edge between the two corresponding neighbors of $v$. By doing this, we have increased the number of connected components by at most one. The result follows. 
    
\end{proof}    

\end{document}
